The paper introduced MeReader, an offline, privacy-preserving, and progress-aware reading companion that constrains both retrieval and generation to a reader's current position in the narrative.
Rather than answering "about the book," it recalls and clarifies without revealing future content or substituting for the act of reading.
Empirically, across three novels and 600 items (300 comprehension, 300 spoiler checks), we observe a consistent quality-latency frontier (Figure~\ref{fig:eval_speed_quality}); \texttt{gemma3n:e4b} and \texttt{llama3.1:8b} achieve the best balance (Bert F1 $\approx 0.74$ -$0.75$ at $\approx 30$s/item, while \texttt{deepseek-r1:1.5b} offers a low-latency alternative ($\approx 15$-$20$s at slightly lower quality (F1 $\approx 0.69$-$0.70$).
Precision-recall distributions (Figure~\ref{fig:eval_precision_recall}) place \texttt{gemma3n:e4b} at the upper frontier (P$\approx0.75$, R$\approx0.74$) with \texttt{llama3.1:8b} and \texttt{gemma3:4b} close behind.
Robustness vs. context availability (Figure~\ref{fig:eval_reading_progress_robustness}) shows stable performance even with $1$-$2$ passages, indicating the system's retrieval pipeline uses evidence that is both compact and sufficient.
Qualitative judging (Figure~\ref{fig:eval_qualitative_metric}) aligns with its quantitative counterparts; the Pareto-efficient models also lead on coherence, contextual fidelity, and helpfulness, i.e., they \textit{feel} better to use, not just score better.

The central novelty is to treat spoiler prevention as a first-class citizen, not an after-the-fact filter.
MeReader's progress gates ensure that generation cannot cite beyond what the reader has already read.
Taken together, the results provide a blueprint for a local, trustworthy reading support on commodity hardware by (1) keeping data and inference on device, (2) incorporating progress metadata into both retrieval and prompting, and (3) targeting frontier-efficient models that deliver strong comprehension at a tolerable latency.
For readers, this design encourages active engagement (asking, checking, connecting) instead of passive consumption.

Our study spans only three (classic) English novels, ten small/medium models when it comes to consumer hardware, and a single workstation configuration.
These may not be enough to form a comprehensive conclusion on said models and the system we used, but are enough to recognise patterns and the possible conclusions we can derive.
The metrics employ BERTScore and an LLM-as-judge, which both carry known biases and do not fully capture long-range discourse reasoning.
The reference answers used, by \texttt{Gemini 2.5 Pro}, were short and may under-reward creative, but at the same time faithful, paraphrasing.
Lastly, the spoiler protocol captures explicit future content but not all forms of implicit foreshadowing.

For future work, we see four directions to improve upon. First, scale evaluation by testing more models on a wider set of prompts/questions and a larger, more diverse pool of books spanning languages and genres; a broader dataset will enable stronger conclusions about MeReader’s effectiveness. Second, broaden generality and accessibility by extending support to multilingual texts and non-fiction. Third, conduct user studies to compare MeReader against baselines, measuring learning outcomes, enjoyment, and engagement. Finally, strengthen evaluation methodology by replacing the judge-as-LLM with human audits.

MeReader shows that progress-aware RAG can deliver spoiler-safe, high-quality reading guidance without leaving the device.
By aligning retrieval, prompting, and evaluation with the reader's position, the system augments, rather than replaces, the act of reading.
We view this as a practical step towards local AI companions for literary engagement and an invitation to the community to refine the blueprint along whatever axes they see fit.

For reviewing the source code of this project and the data used for analysis, please consult the GitHub repository \href{https://github.com/WhiteHades/mereader}{\path{github.com/WhiteHades/mereader}}~\cite{b30}.